\begin{textsample}{POS Dim 1 – rj – Score 31.00 – rj/rj\_ef\_ciencias\_da\_natureza\_1.txt}  \label{ex:f1_pos_279}
3.3.
A Área de Ciências da Natureza A Base Nacional Comum Curricular é um documento de caráter normativo que define o conjunto orgânico e progressivo de aprendizagens essenciais que todos os estudantes devem desenvolver ao longo das etapas e modalidades da Educação Básica.
Nesse contexto, aliada a um regime de colaboração e compromissada com uma educação integral que visa à formação e ao desenvolvimento humano global, foi elaborado o Documento de Orientação Curricular do Estado do Rio de Janeiro de forma alinhada à BNCC.
O ensino de \textbf{Ciências} é um componente fundamental na formação do cidadão \textbf{contemporâneo}.
Segundo Chassot ( 2000 ), a \textbf{Ciência} é considerada “ como uma linguagem para facilitar a leitura do mundo ” e aspira, para além da leitura, a intervenção. “ As necessidades de transformá -lo – e, preferencialmente, transformá -lo em algo melhor ” ( CHASSOT, 2003 ).
A formação de um indivíduo \textbf{crítico} que saiba se posicionar ativamente e intervir na sociedade \textbf{moderna} \textbf{depende} de conhecimentos diversos, tais como éticos-políticos, culturais, filosóficos e, sobretudo, científicos.
Sendo assim, ao longo do Ensino Fundamental, o ensino de \textbf{Ciências} tem um compromisso com o \textbf{letramento} científico.
É o \textbf{letramento} científico um importante exercício da cidadania que facilita a leitura, interpretação e \textbf{compreensão} do mundo e, também, de transformá -lo com base nos aportes \textbf{teóricos} e processuais das \textbf{ciências}.
Sendo assim, o currículo de \textbf{Ciências} deve promover situações que garantam aos estudantes o desenvolvimento das oito \textbf{competências} específicas para o ensino de \textbf{Ciências} da Natureza ( articuladas com as dez \textbf{competências} gerais da Base ) descritas na BNCC ( 2017 ).
No trabalho de construção do DOC-RJ elaborado de \textbf{maneira} alinhada à Base Nacional Comum Curricular, privilegia -se o principal sujeito da ação educativa : o estudante.
Para isso, o professor precisa oportunizar situações de protagonismo, assumindo o seu papel de mediador e apropriando -se de metodologias de aprendizagem ativa.
A metodologia ativa é uma concepção educativa que estimula processos construtivos de ação-reflexão-ação ( FREIRE, 2006 ), em que o estudante participa ativamente de todo processo.
Por exemplo, a Aprendizagem por Resolução de \textbf{Problemas} – ARP ( Metodologia da \textbf{Problematização} ).
Na ARP, os estudantes são estimulados a observar em sua volta e perceberem os \textbf{problemas} que podem ser estudados e solucionados.
De acordo com a BNCC ( 2017 ), o processo investigativo deve ser entendido como elemento-chave na formação dos estudantes.
Para \textbf{tanto}, o aprendizado deve ir além do passo a passo e do cumprimento das etapas pré-definidas do \textbf{método} científico.
É \textbf{indispensável} o estímulo à observação, experimentação e investigação além da \textbf{simples} execução de um roteiro laboratorial, a fim de formar estudantes questionadores e divulgadores do conhecimento científico.
Importante que o docente propicie o diálogo, contextualizando os saberes populares, tradicionais e científicos no ensino de \textbf{Ciências}, contemplando discussões que possam transformar e/ou ressignificar o processo de aprendizagem, respeitando os conhecimentos que fazem parte de nossa diversidade cultural. “ O ensino dialógico-problematizador enriquece ainda mais quando professores desenvolvem suas ações pedagógicas tendo em \textbf{vista} a \textbf{pluralidade} cultural e o amplo espectro de saberes que se acham à sua volta ” ( OLIVEIRA, 2001 ).
Conforme a BNCC ( 2017, p. 329 ), os \textbf{alunos} possuem vivências, saberes, interesses e curiosidades sobre o mundo natural e tecnológico que devem ser valorizados e \textbf{mobilizados}.
Sugestões de práticas \textbf{metodológicas} ativas como pesquisa de campo, práticas experimentais, elaboração de projetos científicos, visitas a centro de \textbf{ciências} e museus, concurso e construção de protótipos, campanhas publicitárias, \textbf{aula} invertida, modelagem, júri simulado, dramatização, \textbf{trilhas} ecológicas, gamificação, feiras científicas, trabalhos em grupo, campanhas solidárias, gincanas interdisciplinares, entre outras, partem do pressuposto da \textbf{contextualização} do cotidiano do estudante.
As metodologias ativas podem ser enriquecidas com a utilização das \textbf{tecnologias} \textbf{digitais}.
No mundo \textbf{contemporâneo}, o \textbf{surgimento} de novas práticas de linguagens abre leques de novas possibilidades de acesso e informação. \textbf{Competências} Específicas de \textbf{Ciências} da Natureza para o Ensino Fundamental 1.
Compreender as Ciências da Natureza como empreendimento humano, e o conhecimento científico como provisório, cultural e histórico. 2.
Compreender conceitos fundamentais e estruturas explicativas das Ciências da Natureza, bem como dominar processos, práticas e procedimentos da investigação científica, de modo a sentir segurança no \textbf{debate} de questões científicas, tecnológicas, \textbf{socioambientais} e do mundo do trabalho, continuar aprendendo e colaborar para a construção de uma sociedade justa, democrática e inclusiva. 3.
Analisar, compreender e explicar características, fenômenos e processos relativos ao mundo natural, social e tecnológico ( incluindo o \textbf{digital} ), como também as relações que se estabelecem entre eles, exercitando a curiosidade para fazer perguntas, buscar respostas e criar soluções ( inclusive tecnológicas ) com base nos conhecimentos das Ciências da Natureza. 4.
Avaliar aplicações e implicações políticas, \textbf{socioambientais} e culturais da \textbf{ciência} e de suas \textbf{tecnologias} para propor alternativas aos desafios do mundo \textbf{contemporâneo}, incluindo aqueles relativos ao mundo do trabalho. 5.
Construir \textbf{argumentos} com base em dados, evidências e informações confiáveis e negociar e defender ideias e pontos de \textbf{vista} que promovam a consciência \textbf{socioambiental} e o respeito a si próprio e ao outro, acolhendo e valorizando a diversidade de indivíduos e de grupos sociais, sem preconceitos de qualquer natureza. 6.
Utilizar diferentes linguagens e \textbf{tecnologias} \textbf{digitais} de informação e comunicação para se comunicar, acessar e disseminar informações, produzir conhecimentos e resolver \textbf{problemas} das Ciências da Natureza de forma \textbf{crítica}, significativa, \textbf{reflexiva} e ética. 7.
Conhecer, apreciar e cuidar de si, do seu corpo e bem-estar, compreendendo -se na diversidade humana, fazendo -se respeitar e respeitando o outro, recorrendo aos conhecimentos das Ciências da Natureza e às suas \textbf{tecnologias}. 8.
Agir pessoal e coletivamente com respeito, autonomia, responsabilidade, flexibilidade, resiliência e determinação, recorrendo aos conhecimentos das Ciências da Natureza para tomar decisões frente a questões científico-tecnológicas e \textbf{socioambientais} e a respeito da saúde individual e coletiva, com base em princípios éticos, democráticos, sustentáveis e solidários.

% matched lemmas: aluno, argumento, aula, ciência, competência, compreensão, contemporâneo, contextualização, crítico, debate, depender, digital, indispensável, letramento, maneira, metodológico, mobilizar, moderno, método, pluralidade, problema, problematização, reflexivo, simples, socioambiental, surgimento, tanto, tecnologia, teórico, trilha, vista
\end{textsample}
